% 附錄 B 結果與 Slurm 工作對照
\chapter{結果與 Slurm 工作對照}\label{app:jobs}

正文中各項結果對應的 Slurm 工作列於下表；原始輸出在 \path{gh200/logs/}，彙整後的結果檔在 \path{gh200/results/}，完整的執行過程（含失敗與重做）記錄在 \texttt{gh200/實驗日誌.md} 與 \texttt{gh200/應用實驗日誌.md}。

{\small\linespread{1.15}\selectfont\renewcommand{\arraystretch}{1.3}
\begin{longtable}{@{}>{\raggedright\arraybackslash}p{6.2cm}>{\raggedright\arraybackslash}p{\dimexpr\textwidth-6.2cm-2\tabcolsep\relax}@{}}
\caption{各結果對應的 Slurm 工作}\label{tab:jobs} \\
\toprule
結果 & Slurm 工作 \\
\midrule
\endfirsthead
\multicolumn{2}{@{}l}{\footnotesize 表~\ref{tab:jobs}（續）} \\
\toprule
結果 & Slurm 工作 \\
\midrule
\endhead
\bottomrule
\endlastfoot
省電機制的填值對照（圖~\ref{fig:mech}） & 1752–1753、1762、1815、1858 \\
矩陣乘法微基準（圖~\ref{fig:gemm}） & 1667–1669 \\
矩陣乘法的誤差模型（第 \ref{sec:errmodel} 節） & 1664 \\
原生 HPL 的調參（第 \ref{sec:hpl-tune} 節） & 3303、3306 \\
調參後 HPL 的功耗上限掃描（表~\ref{tab:hpl-perf}、圖~\ref{fig:hpl}） & 甲組 3315；乙組 3582（3490、3579 因網路問題取消，見第 \ref{sec:hpl-tune} 節） \\
HPL 多組 $N$ 的殘差（表~\ref{tab:hpl-resid}） & 1820（$N = 286720$ 取自 3315、3582） \\
未調參 HPL 的功耗上限掃描 & 1670 \\
HPL-MxP 的功耗上限掃描（圖~\ref{fig:mxp}、表~\ref{tab:mxp-perf}） & 甲組 1671；乙組 3583 \\
HPL-MxP 的功耗煞車診斷、多組 $N$ & 1674、1675；1821、1859 \\
Quantum ESPRESSO 金表面攔截掃描（圖~\ref{fig:qe}） & 1729、1738 \\
科學軟體所需的位數（表~\ref{tab:bits}） & 矽 1728；纈胺黴素 1730、1745（門檻 $10^{-7}$）；紫杉醇 1731 \\
Quantum ESPRESSO 攔截版的效能（表~\ref{tab:qe-perf} 最後一列） & 1737 \\
GPU4PySCF 的剖析（矩陣乘法佔時間的比例） & 1713、1717 \\
以單精度取代 FP64 矩陣乘法（第 \ref{sec:sci-fp32} 節） & 金表面 1928、矽 1930 \\
Quantum ESPRESSO 自動調整（表~\ref{tab:qe-perf}） & 先暖機：金表面 3479、矽 3488；先前未暖機：金表面 1758、2009，矽 1759、2010 \\
Quantum ESPRESSO 不受功耗限制的對照執行 & 2003 \\
方陣上的截斷、模擬與單精度（圖~\ref{fig:square}） & 1894 \\
纈胺黴素真實形狀（圖~\ref{fig:shapes}） & 1894、1896、1897 \\
纈胺黴素矩陣乘法的形狀統計 & 1895 \\
單精度做法的誤差與內積長度（表~\ref{tab:tf32}） & 1934 \\
分界點掃描（圖~\ref{fig:crossover}） & 2005–2008 \\
非矩陣乘法核心（圖~\ref{fig:nongemm}） & 1815、1858、1860 \\
MFEM、SW4、LAMMPS（表~\ref{tab:apps}） & 1992、1995、1996 \\
LoRA 微調（圖~\ref{fig:lora}） & 1732、1736、1760 起、1901、1904、1990、1993；捨入方式 1922、1959；分層截斷 2011、2020、2027 \\
語言模型推論的速度 & 1849 \\
語言模型推論的精度（圖~\ref{fig:llm}、表~\ref{tab:infer}） & 困惑度 1817、1823、1829–1831、1837；注意力 1827、1828；鑰匙檢索 1841；影片 1852、1853；BERT 1857；解碼的省電 1903 \\
DF-MP2（表~\ref{tab:mp2}） & 精度掃描 2037；功耗上限 2038、3489 \\
合成資料的指數跨度（圖~\ref{fig:span}） & 2034–2036 \\
真實應用的跨度與模擬（表~\ref{tab:span-apps}、圖~\ref{fig:rangemap}） & 纈胺黴素 2046、2195；QE 2136、2137；cuBLAS 自動模式 3292；HPL 3304、3305；高斯過程 3248、3291；各應用的逐列跨度統計 2128–2131 \\
寬跨度 DFT 在功耗上限下（第 \ref{sec:gp} 節） & 2821 \\
高斯過程對數概似評估（表~\ref{tab:gp}） & 3248、3501（3482 範圍參數設錯，已取消） \\
高斯過程的舊版截斷實作；小規模 FP32 估計 & 2827；3213 \\
高斯過程最大概似估計、tile 混合精度（表~\ref{tab:gp-mp}） & 3249；3244 \\
ExaGeoStatCPP（表~\ref{tab:exa}） & 3229 \\
\end{longtable}
\par}

MI350X 的結果對應的紀錄列於表~\ref{tab:jobs-mi}。紀錄名稱為「標籤\_日期\_時間」，原始輸出在 \path{mi350/logs/}（節點 n02 的在 \path{mi350/logs/n02/}），彙整後的結果檔在 \path{mi350/results/}，執行過程記錄在 \texttt{mi350/實驗日誌.md}。括號內為 MI350X 叢集的 Slurm 工作編號；沒有編號的是在設定 Slurm 之前直接在節點上執行的。

{\small\linespread{1.15}\selectfont\renewcommand{\arraystretch}{1.3}
\begin{longtable}{@{}>{\raggedright\arraybackslash}p{6.2cm}>{\raggedright\arraybackslash}p{\dimexpr\textwidth-6.2cm-2\tabcolsep\relax}@{}}
\caption{MI350X 各結果對應的紀錄}\label{tab:jobs-mi} \\
\toprule
結果 & 紀錄（Slurm 工作） \\
\midrule
\endfirsthead
\multicolumn{2}{@{}l}{\footnotesize 表~\ref{tab:jobs-mi}（續）} \\
\toprule
結果 & 紀錄（Slurm 工作） \\
\midrule
\endhead
\bottomrule
\endlastfoot
矩陣乘法各上限（圖~\ref{fig:mi-gemm}、表~\ref{tab:mi-robust}） & n01 \texttt{e1\_1028\_1709}；n02 \texttt{e1n02\_1029\_0502}（13） \\
填值對照與鎖頻下的時脈（表~\ref{tab:mi-fill}） & \texttt{e1b\_clk900}（20）、\texttt{e1b\_clk1300}、\texttt{e1b\_clk1300b}（23）、\texttt{e1b\_clk1700}（21）、\texttt{e1c\_lim1300}（24）、\texttt{e1c\_free}（25）；非矩陣乘法核心 \texttt{nongemm}（34） \\
功耗模型（圖~\ref{fig:mi-model}、表~\ref{tab:mi-modelerr}） & 上列矩陣乘法與鎖頻紀錄；外部驗證 \texttt{mp2lock1300}、\texttt{gplock1300}（67） \\
HPL 的安全位數（表~\ref{tab:mi-hplres}） & \texttt{hplbits262k}、\texttt{hplbits393k}、\texttt{hplbits540k}（17–19） \\
HPL 各上限 & n01 \texttt{hpl\_1029\_0336}；n02 \texttt{hpln02\_1029\_1045}（15） \\
HPL-MxP 各上限 & n01 \texttt{mxp\_1028\_2359}；n02 \texttt{mxpn02\_1029\_0810}（14）；800~W 成對重跑 \texttt{mxp800pair}（50）；精確對角分解 \texttt{mxpexact}（65） \\
DF-MP2 & 張量準備 \texttt{mp2prep}（31、58）與驗證 \texttt{mp2val}（29、30）；n01 \texttt{mp2\_1029\_1607}（32）；n02 \texttt{mp2n02}（59）；GEMMul8 \texttt{mp2emu\_1029\_1741}（40） \\
高斯過程 & n01 \texttt{gp\_1029\_0731}；n02 \texttt{gpn02\_1030\_0521}（57）；GEMMul8 $n=80000$ \texttt{gpemu}（52）、$n=64000$ \texttt{gpemu64k}（63） \\
截斷與 GEMMul8 的分界（圖~\ref{fig:mi-xover}） & \texttt{xover\_1029\_1508}（28） \\
HPL 與 GEMMul8（表~\ref{tab:mi-h2h}） & 最佳設定 \texttt{hpl3B}（49）；相近精度 \texttt{hpl3B2}（64）；$1\times8$ 工作空間不足 \texttt{hpl3A}（43、48）、\texttt{hpl3A2}（54） \\
GEMMul8 在 $mn=2^{32}$ 的錯誤 & \texttt{g8thr}（66） \\
指數範圍、模擬內再截斷 & \texttt{range\_*}（53）；\texttt{synergy}（51） \\
穩態與熱預處理（第 \ref{sec:mi-temp} 節） & \texttt{steady\_1030\_0901}（68）；\texttt{mxpprecond\_1030\_0944}（69） \\
位元變化對照（表~\ref{tab:mi-flips}） & \texttt{mech1300}（102）、\texttt{mechb1300}（109）、\texttt{mechb1700}（110） \\
高斯過程的三參數最大概似估計 & \texttt{gpmle}（111） \\
DF-MP2 每次呼叫的時間 & \texttt{mp2call}（112） \\
GEMMul8 3.5.4 的溢位與初始化成本 & \texttt{g8thr354}（122）、\texttt{g8cold}（123） \\
相同精度的比較（圖~\ref{fig:mi-pareto}） & \texttt{pareto}（104） \\
逐次呼叫選擇的形狀掃描（表~\ref{tab:mi-selector}） & \texttt{shpA}（159）、\texttt{shpB}（160） \\
\end{longtable}
\par}

